\chapter{評価}

本章では，第3章で提案・実装した災害分析基盤を約9ヶ月間（2025年2月中旬〜10月）運用し，その有効性を検証した結果を述べる．
まず運用実績とシステムの安定性を示し，次に収集データの統計的特性を分析する．
最後に，地震発生時のコメント反応や感情分析結果を通じて，実際に定性的なYouTubeライブ配信コメントを定量的なデータとして扱いどのような特徴があるかを明らかにする．
\section{運用実績とシステム検証}

\subsection{データ収集実績}

2025年2月15日〜10月31日の約8.5ヶ月間で，蓄積したデータについて評価する。
\tabref{tab:key_fields_youtube_live_comments}に示したsnippet\_typeフィールドの種類別に集計した結果を\tabref{tab:data_collection_results}に示す．

\begin{table}[h]
    \centering
    % \caption{YouTube収集データの種類別集計}
    \caption{YouTube collection data summary by type}
    \label{tab:data_collection_results}
    \begin{tabular}{lrr}
        \hline
        イベント種別 (snippet\_type) & 件数     & 割合 (\%)  \\
        \hline
        \hline
        textMessageEvent       & 472996  & 97.78    \\
        userBannedEvent        & 10165   & 2.10     \\
        superChatEvent         & 527    & 0.11     \\
        superStickerEvent      & 48     & 0.01(未満) \\
        pollEvent              & 2      & 0.01(未満) \\
        sponsorOnlyModeEndedEvent & 4   & 0.01(未満) \\
        \hline
        合計                     & 483742 & 100.00   \\
        \hline
    \end{tabular}
\end{table}

本来取得が可能であった全体のコメントの数を把握できないため、以下にシステム運用中に見つかったエラーやデータ欠損の要因をし、収集したデータ全体からコメントデータが十分に確保できているかを検証した．

まず，運用中の2025年7月6日早朝から半日程度，検証用の別システムの影響により，システムが停止しデータが欠損ていたことが確認された．

また，収集対象のライブ配信では地震規模が一定レベルを超えた場合，書込み制限が起こるため，システムの停止以外に最低でも5分程度欠損が起こる場合が確認された．
その上で，YouTubeでは，1日にリクエストできる上限と，1度のリクエスト含まれるコメントの上限は2000個である．
しかし，実際に一度のリクエストに含まれたコメントは最高で397件であったため，システムの稼働中はデータが確実に収集できていると考えられる．

収集した全483742件のデータのうち，textMessageEventが大半を占めており，分析対象となるコメントデータは全体の約98\%を占めている．
システムの停止期間は全体に比べて小さくコメントデータは十分に確保できていると考えられる．

%TODO: 
% \subsection{API利用状況とQuota検証}

% 実際の運用で以下を確認した．

% \begin{itemize}
%     \item \textbf{Quota消費量の実測}: GCPダッシュボードで確認した結果，1リクエストあたり5 Quotaを消費することを実証．公式ドキュメントに記載がないため，今後の仕様変更に注意が必要．
%     \item \textbf{2分間隔の妥当性}: 9ヶ月間の運用で1日のQuota上限（10,000）に到達せず，平均消費量は約3,600 Quota/日（720リクエスト）であった．
%     \item \textbf{取得漏れの検証}: 地震発生後の高頻度投稿時も，2分間隔で十分にコメントを捕捉できることを確認．
% \end{itemize}

